\section{Empirical Evaluation}
\label{sec:synthetic}

\subsection{Experimental Overview}
\label{sec:empirical-overview}

The empirical evaluation addresses three questions that follow from the calibration and design results in Sections 4--5. First, with only finitely many dense calibration trajectories, how accurately can protocol values be estimated across a candidate family, and at what resolution can competing protocols be reliably distinguished and selected? Second, once those values are estimable, does designing measurements for the predictive value of the specified target improve on target-free or target-mismatched objectives, and how closely can practical search approach the exact target-specific optimum? Third, in fully annotated temporal data, are broad matched-budget differences in observation layout reproducible, and is there sufficient information to support the finer, data-dependent choice of exact observation locations?


The first two questions are studied under known data-generating laws. The
calibration experiments compare plug-in values with the exact population values
$\cI_g(S;K)$ and evaluate family-uniform error, selection regret and nested
candidate classes. The design experiments enumerate every feasible protocol in
moderate catalogues, so the target-specific optimum and the optimisation error
of forward selection with one-swap refinement are known exactly. The third
question is studied retrospectively in Sleep-EDF and Long-Term AF by
reconstructing partial observation protocols from complete annotations. Because
population protocol values are not observed in these data, performance is
measured by pooled held-out cross-fitted $R^2$; support selection and predictor
tuning are confined to the corresponding outer training folds. Simulation
specifications are given in Appendix~\ref{sec:appendix-simulations}, and the
real-data processing, cross-fitting and resampling procedures are given in
Appendix~\ref{sec:appendix-real}.



\subsection{Finite Calibration and Protocol Comparison}
\label{sec:simulation-evidence}
\label{sec:calibration-regret}

\Cref{fig:calibration} examines three consequences of finite calibration.
Panel (a) asks how accurately protocol values can be estimated over an entire
candidate family. Panel (b) asks whether the remaining estimation error is
large enough to affect selection of a near-optimal protocol. Panels (c,d)
then examine how calibration accuracy limits the useful granularity of the
candidate family. In all experiments, protocol values estimated from the
calibrated covariance $\widehat K$ are compared with their exact population
values under the known generating covariance. Full simulation specifications
are given in Appendix~D.

\begin{figure}[H]
\centering
\includegraphics[width=0.82\linewidth]{figures/fig_calibration.pdf}
\caption{Finite-calibration error, selection regret and candidate-family
resolution.
(a) Family-uniform protocol-value error $\varepsilon_m$ for the temporal-mean
and occupation-at-zero targets over all \numEstFamily{} four-action protocols
under a stationary OU Gaussian model on a $p=128$ grid.
(b) True regret of the protocol selected from estimated values, together with
the $2\varepsilon_m$ envelope from \cref{sec:estimation}, for four-action
selection from a 14-action catalogue, averaged over two covariance models and
three targets: the temporal mean and occupation above thresholds zero and one.
(c,d) Selection regret and family-uniform error for four nested candidate
families of increasing granularity, with sizes
\numResClassSizeOne, \numResClassSizeTwo, \numResClassSizeThree{} and
\numResClassSizeFour, under a nonstationary $p=64$ covariance.
Calibration sizes are
$m\in\{25,50,100,250,500,1000\}$.
Panels (a,b) use \numEstReplications{} replications and panels (c,d) use 30;
curves show replication means and bands show standard errors.}
\label{fig:calibration}
\end{figure}

Panel (a) evaluates the family-uniform error
$\varepsilon_m=\max_{S\in\Pi_{\mathcal B}}
|\widehat{\cI}_g(S)-\cI_g(S;K)|$,
the largest protocol-value error within the candidate family.
This error decreases steadily as the number of calibration trajectories
increases. The full-range log--log slopes are \numSlopeMean{} and
\numSlopeOccZero; over the three largest calibration sizes they are
\numTailSlopeMean{} and \numTailSlopeOccZero, approaching the fixed-model
exponent $-1/2$ derived in \cref{sec:estimation}.

In this experiment, panel (b) shows that selecting a near-optimal protocol can
require substantially less accuracy than estimating every protocol value
uniformly. Across the six covariance--target cells, the mean cellwise log--log
slopes are \numRegretSlope{} for realised selection regret and
\numEpsSlope{} for the uniform-error envelope. The median and maximum
ratios of regret to this envelope are \numRegretRatioMedian{} and
\numRegretRatioMax. The difference reflects the two criteria being controlled:
$\varepsilon_m$ is determined by the largest estimation error anywhere in the
family, whereas near-optimal selection depends primarily on correctly ordering
the leading candidates.

Panels (c,d) show how this distinction affects the useful resolution of the
candidate family. At $m=25$, the finest class has selection regret
\numResRegretFineSmall, compared with \numResRegretCoarseSmall{} for the
coarsest class, despite a family-uniform error of \numResEpsFourSmall.
Thus, a richer candidate class can yield a better selected protocol even
before all of its values are estimated accurately. By $m=1000$, the uniform
error of the finest class falls to \numResEpsFourLarge, whereas the coarsest
class still incurs restriction regret \numResRegretCoarseLarge{} because
better protocols lie outside that class.

Together, the four panels distinguish uniform estimation resolution from
selection resolution. Reliable selection does not require every candidate
value to be estimated to the same precision; it requires sufficient
calibration to resolve the value differences among the candidates that compete
near the optimum. Candidate-family refinement can therefore be useful before
uniform estimation over the refined family becomes accurate.



\subsection{Target-Aware Design and Search}
\label{sec:target-design-experiment}

\Cref{tab:s5} examines two questions about observation design. First, how much
is lost when protocols are chosen using target-free or target-mismatched
objectives rather than the predictive value of the specified target? Second,
when the target-aware objective is used, how closely does the proposed search
procedure approach the exact optimum?

We evaluate \numDesignComparableInstances{} instances formed by five aggregate
targets and five covariance--action settings: stationary OU, stationary
Mat\'ern, horizon-varying covariance, recency-weighted prediction and
heterogeneous actions. The first four settings select four of 12 point actions;
the heterogeneous setting chooses among 18 actions under a cost budget of four.
Every feasible protocol is enumerated, providing the exact target-specific
optimum against which both alternative design objectives and approximate search
can be evaluated. Full simulation settings are given in Appendix~D.

For a method returning protocol $S$, we report relative efficiency
$\cI_g(S)/\cI_g(S^*)$, where $S^*$ is the exhaustive optimum for the target
being evaluated. An efficiency of one therefore indicates that the selected
protocol attains the target-specific optimum. Latent-state mutual information
and integrated posterior variance provide target-free comparators. The
linear-target comparator retains the actual temporal weights and measurement
noise but replaces the target transformation by $g(z)=z$, whereas kernel
quadrature uses the corresponding conventional noiseless integration
criterion.

\begin{table}[H]
\centering\small
\begin{tabular}{@{}lrrr@{}}
\toprule
Method & minimum & mean & median \\
\midrule
Latent-state mutual information & $\numDesignMethodMinMI$ & $\numDesignMethodMeanMI$ & $\numDesignMethodMedianMI$ \\
Integrated posterior variance & $\numDesignMethodMinIMSE$ & $\numDesignMethodMeanIMSE$ & $\numDesignMethodMedianIMSE$ \\
Actual-noise linear target & $\numDesignMethodMinLinear$ & $\numDesignMethodMeanLinear$ & $\numDesignMethodMedianLinear$ \\
Noiseless kernel quadrature & $\numDesignMethodMinKQ$ & $\numDesignMethodMeanKQ$ & $\numDesignMethodMedianKQ$ \\
Target-aware greedy & $\numDesignMethodMinGreedy$ & $\numDesignMethodMeanGreedy$ & $\numDesignMethodMedianGreedy$ \\
Target-aware greedy plus one-swap & $\numDesignMethodMinSwap$ & $\numDesignMethodMeanSwap$ & $\numDesignMethodMedianSwap$ \\
\bottomrule
\end{tabular}
\caption{Relative efficiency of alternative design objectives and search
procedures against the exhaustive target-specific optimum over
\numDesignComparableInstances{} kernel--action--target instances.}
\label{tab:s5}
\end{table}
\FloatBarrier

The choice of design objective matters most when temporal symmetry is broken.
In the stationary, uniformly weighted settings, supports optimised for
different targets transfer with little loss. Under nonstationarity,
non-uniform temporal weighting or heterogeneous actions, the target-specific
optimal supports become more distinct; the lowest cross-target efficiency in
the heterogeneous setting is \numDesignCrossHetero. The target-free
comparators show the same general distinction: good coverage of the latent
trajectory does not necessarily yield a protocol that is optimal for the
specified aggregate target.

The proposed search procedure is nevertheless close to the exhaustive
target-aware optimum. One-swap refinement raises the minimum efficiency in the
heterogeneous setting from \numDesignGreedyMinHetero{} to
\numDesignAwareMinHetero{} and reaches the exhaustive optimum there and in the
stationary OU and horizon-varying settings. In the remaining settings its
minimum efficiencies are \numDesignAwareMinMatern{} and
\numDesignAwareMinRecency, while evaluating only
\numDesignEvalFracHetero--\numDesignEvalFracStat{} of the exhaustive sets.












\subsection{Retrospective Protocol Comparison}
\label{sec:real}

The retrospective analyses use two fully annotated temporal data sets to
reconstruct partial observation protocols while computing the prediction target
from the complete analysable record. 

Sleep-EDF Expanded
\citep{kemp2000analysis,mourtazaev1995age,goldberger2000physiobank,
pollard2026physionet} contains 30-second sleep-stage annotations from the
Sleep Cassette (SC) and Sleep Telemetry (ST) studies. The analysed sample
contains \numSleepRecordings{} recordings from \numSleepSubjects{} subjects and
\numSleepHours{} hours of valid annotations. Stages 3 and 4 are combined as N3
\citep{rechtschaffen1968manual,iber2007aasm}. Each recording is mapped to
$p=\numSleepGrid$ relative-time anchors. The targets are the exact full-interval
REM, N3 and Wake proportions. The Long-Term AF Database
\citep{petrutiu2007abrupt,goldberger2000physiobank} contains
\numLtafRecords{} records and \numLtafHours{} hours of reviewed rhythm
annotations, with median analysable record length \numLtafMedianHours{} hours.
Analysis begins at the first rhythm marker, excluding any unannotated prefix.
Each record is mapped to $p=\numAfGrid$ equal relative-time bins, and the target
is the exact fraction of analysable annotated time spent in atrial fibrillation.

Sleep protocols expose stage-specific binary indicators at selected
relative-time anchors, with budgets $N\in\{4,8,16,32,64\}$ counting distinct
scored 30-second epochs. AF protocols expose bin-level AF fractions, with
budgets $N\in\{1,2,4,8,16,32\}$ observing the fraction $N/p$ of each analysable
record. Both analyses use five outer folds and report $R^2$ from pooled held-out
predictions. Sleep folds are subject-disjoint and stratified by source study,
with subject-balanced evaluation weights. Study, valid-duration and treatment
covariates enter the Sleep predictors unpenalised. AF uses record-level folds
and equal record weights because repeated-subject identifiers are unavailable.

At each budget, we compare a centred contiguous block with observations
dispersed as uniformly as possible over relative time. Sleep additionally
learns target-aware and kernel-quadrature supports within each outer training
fold. The target-aware criterion maximises a regularised best-linear estimate
of the fraction of residual target variance explained by the selected anchors.
The kernel-quadrature comparator uses the same repaired temporal covariance
without target outcomes. Support selection, ridge tuning and predictor fitting
are completed within the outer training fold. For fixed-template contrasts, conditional 2.5th--97.5th percentile ranges are
obtained by resampling held-out subjects for Sleep and records for AF while
keeping the fitted predictors fixed. Further details of the common-grid
mapping, weighting, support selection and resampling are given in
Appendix~\ref{sec:appendix-real}.

For REM, \cref{fig:real}(a) shows cross-fitted $R^2$ values of
\numCfContigFour, \numCfContigSixteen{} and \numCfContigSixtyFour{} for the
centred-contiguous template at $N=4,16,64$, compared with
\numCfUniformFour, \numCfUniformSixteen{} and \numCfUniformSixtyFour{} for
uniform dispersion. Across the 15 Sleep target--budget cells, the
dispersed-minus-contiguous percentile range has a positive lower endpoint in
\numSleepAdjustedPositiveRanges{} cells. The pooled Sleep results favour dispersion at several budgets, but the
source-specific contrasts are heterogeneous. In separate SC and ST analyses,
\numSleepCohortPositiveRanges{} of \numSleepCohortCells{}
source--target--budget ranges lie above zero,
\numSleepCohortNegativeRanges{} lie below zero and
\numSleepCohortUnresolvedRanges{} include zero. Additional analyses are
reported in Appendix~\ref{sec:appendix-real-sensitivity}.

For AF, at $N=4$ each template observes \numAfObservedPctFour{} of the
analysable record, equivalent to \numAfEquivalentHoursFour{} hour on a 24-hour
record. The contiguous and dispersed templates attain cross-fitted $R^2$
values \numCfAfContigOneHour{} and \numCfAfDispOneHour, respectively, with
difference \numAfBootDiffFour{} and conditional paired percentile range
\numAfBootRangeFour. At $N=16$, the corresponding values are
\numCfAfContigFourHour{} and \numCfAfDispFourHour, with difference
\numAfBootDiffSixteen{} and range \numAfBootRangeSixteen. The paired percentile range has a positive lower endpoint at every AF budget
containing at least two windows. At $N=1$, the difference is
\numAfBootDiffOne{} with range \numAfBootRangeOne. Because only one window is
observed, this comparison reflects window location rather than dispersion.
Outcome-defined cohort and grid-discretisation analyses are reported in
Appendix~\ref{sec:appendix-real-sensitivity}.

\FloatBarrier
\begin{figure}[htbp]
\centering
\includegraphics[width=0.86\linewidth]{figures/fig_real.pdf}
\caption{Held-out prediction from reconstructed Sleep and AF observation
protocols. (a) Pooled baseline-adjusted cross-fitted $R^2$ for the Sleep
targets under target-aware, learned kernel-quadrature, uniformly dispersed and
contiguous supports. (b) Pooled cross-fitted $R^2$ for AF burden under
dispersed and contiguous window layouts.}
\label{fig:real}
\end{figure}
\FloatBarrier

We next examine the stability of the learned Sleep supports. At $N=16$, the
target-aware REM support attains baseline-adjusted cross-fitted
$R^2=\numCfAwareSixteen$, compared with \numCfKqSixteen{} for learned kernel
quadrature and \numCfUniformSixteen{} for fixed dispersion. The fold-averaged
pairwise Jaccard overlap among the REM-, N3- and Wake-specific supports is
\numSleepJaccardMin--\numSleepJaccardMax. The fitted supports therefore differ
substantially across targets, although this difference may also reflect
finite-sample variation in support selection.

\Cref{fig:sweep}(a) varies the number of subjects used for support selection
from \numSweepMSmall{} to the complete outer-training-fold sample while keeping
predictor fitting on the full training fold. The target-aware difference changes
from \numSweepDeltaKqSmall{} to \numSweepDeltaKqFull{} relative to learned
kernel quadrature and from \numSweepDeltaUniSmall{} to
\numSweepDeltaUniFull{} relative to fixed dispersion. Panel (b) repeats the complete selection and evaluation pipeline on
\numSelectionSubsampleReps{} study-stratified subsamples containing
\numSelectionSubsamplePct{} of subjects. The target-aware difference relative
to kernel quadrature has median \numSelectionDeltaKqMed{} and 2.5th--97.5th
percentile range
$[\numSelectionDeltaKqPlo,\numSelectionDeltaKqPhi]$. Relative to fixed
dispersion, the corresponding values are \numSelectionDeltaUniMed{} and
$[\numSelectionDeltaUniPlo,\numSelectionDeltaUniPhi]$. Both ranges include
zero.

\begin{figure}[htbp]
\centering
\includegraphics[width=0.84\linewidth]{figures/fig_sweep.pdf}
\caption{Stability of the learned REM support at $N=16$.
(a) Target-aware held-out advantage as the support-selection sample increases.
(b) Target-aware advantage over learned kernel quadrature and fixed dispersion
across \numSelectionSubsampleReps{} study-stratified
\numSelectionSubsamplePct{} subject subsamples.}
\label{fig:sweep}
\end{figure}
\FloatBarrier

The real-data results distinguish broad temporal-layout contrasts more clearly
than exact learned support locations. Dispersion performs better than contiguous
observation in several pooled Sleep comparisons and across the multi-window AF
budgets, although the Sleep contrasts vary between source studies. The learned
Sleep supports are less stable across targets and subject subsamples, and their
paired held-out $R^2$ differences include zero. This agrees with the resolution
pattern in \cref{fig:calibration}: coarse protocol differences can be easier to
distinguish than fine support locations with finite data.

